\documentclass[tikz,border=4pt]{standalone}
\usepackage{ctex}
\usepackage{amsmath,amssymb}
\usepackage{pgfplots}
\pgfplotsset{compat=1.18}
\usepgfplotslibrary{groupplots}
\begin{document}
\begin{tikzpicture}
\begin{groupplot}[
  group style={
    group size=2 by 1,
    horizontal sep=1.8cm,
  },
  width=6.5cm, height=5.5cm,
  xlabel={$x$}, ylabel={$y$},
  xlabel style={font=\small},
  ylabel style={font=\small},
  tick label style={font=\footnotesize},
  grid=both, grid style={gray!20},
  axis lines=center,
  samples=200,
  legend style={font=\footnotesize, at={(0.98,0.98)}, anchor=north east},
]

% Plot 1: Convex function x^2 with chord above
\nextgroupplot[
  title={凸函数：弦线在曲线上方},
  title style={font=\small},
  xmin=-2.5, xmax=2.5,
  ymin=-0.5, ymax=6.5,
]
\addplot[blue, very thick, domain=-2.2:2.2] {x^2};
\addlegendentry{$f(x)=x^2$（凸）}
% chord from (-2,4) to (2,4)
\addplot[red, dashed, thick] coordinates {(-2,4)(2,4)};
\addlegendentry{弦线 $\geq f$}
% tangent at x=0.8
\addplot[orange, thick, domain=-1.5:2.5] {1.6*x - 0.64};
\addlegendentry{切线 $\leq f$}
% points
\addplot[mark=*, mark size=2.5pt, red] coordinates {(-2,4)(2,4)};
\node[font=\footnotesize, above right] at (axis cs:0.2, 0.5) {$f(\lambda x_1\!+\!(1\!-\!\lambda)x_2)$};
\node[font=\footnotesize, above right] at (axis cs:0.2, 3.0) {$\leq\lambda f(x_1)\!+\!(1\!-\!\lambda)f(x_2)$};

% Plot 2: Nonconvex function sin
\nextgroupplot[
  title={非凸函数：弦线可穿越曲线},
  title style={font=\small},
  xmin=-0.5, xmax=7,
  ymin=-1.8, ymax=2.2,
]
\addplot[blue, very thick, domain=0:6.5] {sin(deg(x))};
\addlegendentry{$f(x)=\sin x$（非凸）}
% chord from (0.5, sin(0.5)) to (5.5, sin(5.5))
\addplot[red, dashed, thick] coordinates {(0.5, 0.479) (5.5, -0.706)};
\addlegendentry{弦线可在曲线下方}
\addplot[mark=*, mark size=2.5pt, red] coordinates {(0.5,0.479)(5.5,-0.706)};
\node[font=\footnotesize, fill=white] at (axis cs:3.5, 0.9) {局部极值 $\neq$ 全局最优};
% mark local max
\addplot[mark=triangle*, mark size=3pt, green!60!black] coordinates {(1.571, 1)};
\node[font=\footnotesize, above] at (axis cs:1.571, 1) {局部};
\addplot[mark=triangle*, mark size=3pt, orange] coordinates {(4.712, -1)};
\node[font=\footnotesize, below] at (axis cs:4.712, -1) {局部};

\end{groupplot}

\node[font=\normalsize\bfseries] at (8, 6.2) {凸函数 vs 非凸函数：几何直觉对比};
\end{tikzpicture}
\end{document}
